\paragrafos{Απόσταση σημείου από ευθεία και απόσταση παραλλήλων}{Β.1.10, «Απόσταση σημείου από ευθεία – Απόσταση παραλλήλων»}

\begin{anagke}
Στον κήπο υπάρχει ένα ευθύγραμμο μονοπάτι και ένα δέντρο που βρίσκεται λίγο πιο πέρα, εκτός μονοπατιού. Ο Νίκος θέλει να φτάσει από το δέντρο στο μονοπάτι όσο πιο γρήγορα γίνεται. Δοκιμάζει τρεις δρόμους με ένα σχοινί: έναν λοξό προς τα αριστερά, έναν λοξό προς τα δεξιά και έναν κατευθείαν «πάνω». Ο ευθύς είναι ο μικρότερος. «Ο δρόμος που πέφτει κάθετα στο μονοπάτι είναι ο συντομότερος!» ανακαλύπτει.

Η Αγγελική προσθέτει ένα δεύτερο σχέδιο: δύο παράλληλες σειρές με δέντρα. «Πόσο απέχουν οι δύο σειρές;» Μετράει από ένα δέντρο της μιας στην άλλη, κάθετα, και βρίσκει $2{,}4$ μέτρα. Μετράει από ένα άλλο δέντρο και βρίσκει πάλι $2{,}4$. «Σε όλα τα σημεία η απόσταση είναι ίδια!». Η Πολυξένη το συμπληρώνει: «Γι' αυτό οι ράγες του τρένου έχουν σταθερή απόσταση: αλλιώς το τρένο δεν χωράει».

Σε αυτή την παράγραφο ορίζουμε την απόσταση ενός σημείου από μια ευθεία και την απόσταση δύο παραλλήλων, δύο έννοιες που θα τις χρειαζόμαστε συνεχώς.
\end{anagke}

\begin{lexi}{απόσταση, κάθετο, πλάγιο}{απόσταση, κάθετο, πλάγιο}
Η \lexiword{απόσταση} σημείου από ευθεία είναι ο «συντομότερος δρόμος» από το σημείο προς την ευθεία. Το \lexiword{κάθετο} τμήμα πέφτει με ορθή γωνία πάνω στην ευθεία, ενώ τα υπόλοιπα, τα \lexiword{πλάγια}, είναι «λοξά» και μακρύτερα.
\end{lexi}

\begin{metafrasi}
Το κάθετο τμήμα από το $\mathrm{Α}$ στην $(\varepsilon)$. & $\mathrm{ΑΜ}\perp(\varepsilon)$\\
Η απόσταση του $\mathrm{Α}$ από την $(\varepsilon)$. & $\mathrm{ΑΜ}$\\
Η απόσταση δύο παραλλήλων. & η απόσταση οποιουδήποτε σημείου της μιας από την άλλη\\
Κάθε πλάγιο τμήμα. & μεγαλύτερο από το κάθετο\\
\end{metafrasi}

\tmima{Η θεωρία}

\begin{orismos}{απόσταση σημείου από ευθεία}
\textbf{Απόσταση} ενός σημείου $\mathrm{Α}$ από μια ευθεία $(\varepsilon)$ είναι το μήκος του \textbf{κάθετου} τμήματος $\mathrm{ΑΜ}$ από το $\mathrm{Α}$ προς την $(\varepsilon)$, όπου $\mathrm{Μ}$ το ίχνος της καθέτου.
\end{orismos}

Το κάθετο τμήμα είναι το \textbf{συντομότερο} από όλα τα τμήματα που ενώνουν το $\mathrm{Α}$ με σημεία της $(\varepsilon)$. Κάθε \textbf{πλάγιο} τμήμα είναι μεγαλύτερο από το κάθετο. Αν το σημείο ανήκει στην ευθεία, η απόστασή του από αυτήν είναι $0$.

\begin{center}
\begin{tikzpicture}[scale=1.1]
\draw[thick,{Stealth}-{Stealth}] (-2.5,0)--(3.5,0) node[right] {$(\varepsilon)$};
\fill (0.5,2) circle (2pt) node[above] {$\mathrm{Α}$}; \fill (0.5,0) circle (2pt) node[below] {$\mathrm{Μ}$};
\draw[gblue,thick] (0.5,2)--(0.5,0); \draw (0.5,0.25)--(0.75,0.25)--(0.75,0);
\draw (0.5,2)--(-1.5,0); \draw (0.5,2)--(2.5,0);
\node[gblue,right] at (0.5,1) {$\mathrm{ΑΜ}$};
\end{tikzpicture}
\end{center}

\tmimasmall{Απόσταση παραλλήλων}

Αν $\varepsilon_1\parallel\varepsilon_2$, όλα τα σημεία της $\varepsilon_1$ απέχουν την ίδια απόσταση από την $\varepsilon_2$. Η απόσταση αυτή λέγεται \textbf{απόσταση των παραλλήλων}. Δύο παράλληλες ευθείες έχουν σταθερή απόσταση, δηλαδή δεν πλησιάζουν ούτε απομακρύνονται.

Το σύνολο των σημείων που απέχουν απόσταση $\delta$ από μια ευθεία $(\varepsilon)$ είναι δύο παράλληλες ευθείες προς την $(\varepsilon)$, μία σε κάθε ημιεπίπεδο, που απέχουν $\delta$ από αυτήν.

\begin{prosoxi}
\begin{itemize}[leftmargin=1.2em,itemsep=0.3ex,topsep=2pt]
\item Η απόσταση σημείου από ευθεία μετριέται πάντα κάθετα, όχι λοξά.
\item Το πλάγιο τμήμα δεν δείχνει την απόσταση: είναι μεγαλύτερο.
\item Η απόσταση παραλλήλων είναι η ίδια όπου κι αν τη μετρήσεις, αλλά πρέπει να τη μετράς κάθετα.
\item Η απόσταση ενός σημείου της ευθείας από την ευθεία είναι $0$.
\end{itemize}
\end{prosoxi}

\tmima{Λυμένα παραδείγματα}

\paradeigma{Ο συντομότερος δρόμος}
\tixero Ένα δέντρο $\mathrm{Α}$ εκτός ευθύγραμμου μονοπατιού $(\varepsilon)$.\par
\tithelo Ποιος δρόμος από το δέντρο προς το μονοπάτι είναι ο συντομότερος.\par
\possindeo Η απόσταση του $\mathrm{Α}$ από την $(\varepsilon)$ είναι το κάθετο τμήμα $\mathrm{ΑΜ}$. Κάθε πλάγιο $\mathrm{ΑΒ}$ είναι μεγαλύτερο.\par
\apantisi Ο κάθετος δρόμος $\mathrm{ΑΜ}$.\par
\elegxos Με σχοινί: ο λοξός δρόμος βγαίνει μακρύτερος.

\paradeigma{Η απόσταση παραλλήλων}
\tixero Δύο σειρές δέντρων παράλληλες και $\mathrm{ΑΜ}\perp\varepsilon_2$ με $\mathrm{ΑΜ}=2{,}4$ μ.\par
\tithelo Η απόσταση ενός άλλου δέντρου $\mathrm{Β}$ της $\varepsilon_1$ από την $\varepsilon_2$.\par
\possindeo Οι παράλληλες έχουν σταθερή απόσταση.\par
\apantisi $2{,}4$ μ.\par
\elegxos Κάθε κάθετο τμήμα ανάμεσα στις παράλληλες είναι $2{,}4$ μ.

\paradeigma{Τρεις σειρές}
\tixero Τρεις παράλληλες σειρές δέντρων σε ίσες αποστάσεις $2{,}4$ μ.\par
\tithelo Η απόσταση της πρώτης από την τρίτη.\par
\possindeo Η απόσταση από την πρώτη στη δεύτερη είναι $2{,}4$ και από τη δεύτερη στην τρίτη $2{,}4$, αφού οι παράλληλες είναι διαδοχικές. $2{,}4+2{,}4=4{,}8$.\par
\apantisi $4{,}8$ μ.\par
\elegxos Η μεσαία είναι στη μέση των δύο άλλων.

\paradeigma{Σημεία με δοσμένη απόσταση}
\tixero Μια ευθεία $(\varepsilon)$.\par
\tithelo Ποια σημεία απέχουν $3$ cm από αυτήν.\par
\possindeo Σε κάθε ημιεπίπεδο υπάρχει μία παράλληλη προς την $(\varepsilon)$ σε απόσταση $3$ cm.\par
\apantisi Δύο παράλληλες ευθείες, μία σε κάθε ημιεπίπεδο.\par
\elegxos Η απόσταση ανάμεσά τους είναι $6$ cm.

\tmima{Ασκήσεις}

\zoni{Το κατάλαβα}{κάθετα και πλάγια}

\ex Σχεδίασε μια ευθεία $(\varepsilon)$ και ένα σημείο $\mathrm{Α}$ εκτός της. Φέρε την κάθετη από το $\mathrm{Α}$ και δύο πλάγια. Μέτρησε και σύγκρινε τα μήκη τους.

\ex Μέτρησε την απόσταση του $\mathrm{Α}$ από την $(\varepsilon)$. Πόση είναι η απόσταση ενός σημείου που βρίσκεται πάνω στην $(\varepsilon)$;

\ex Σχεδίασε δύο παράλληλες ευθείες και μέτρησε την απόστασή τους σε τρία διαφορετικά σημεία. Τι παρατηρείς;

\ex Τρεις παράλληλες ευθείες απέχουν διαδοχικά $2{,}4$ cm. Πόση είναι η απόσταση της πρώτης από την τρίτη;

\ex Σχεδίασε όλα τα σημεία που απέχουν $2$ cm από μια ευθεία $(\varepsilon)$.

\ex Σωστό ή λάθος; α) Το κάθετο τμήμα είναι το μικρότερο από όλα όσα ενώνουν ένα σημείο με την ευθεία. \quad β) Η απόσταση παραλλήλων δεν είναι σταθερή. \quad γ) Η απόσταση σημείου της ευθείας από την ευθεία είναι $0$. \quad δ) Το πλάγιο είναι μικρότερο από το κάθετο.

\zoni{Το χρησιμοποιώ}{κήπος, δρόμοι, σχοινιά}

\ex \textbf{Το κοντινότερο.} Ένα δέντρο απέχει $3$ μ. κάθετα από ένα μονοπάτι. Ο Νίκος δοκιμάζει να φτάσει λοξά και βρίσκει $3{,}6$ μ. Μπορεί να έχει υπολογίσει τον κάθετο δρόμο λάθος; Εξήγησε.

\ex \textbf{Οι σειρές.} Δύο παράλληλες σειρές λαχανικών απέχουν $60$ cm. Μια τρίτη σειρά, παράλληλη προς αυτές και προς την άλλη πλευρά της δεύτερης, απέχει $60$ cm από τη δεύτερη. Πόσο απέχει από την πρώτη;

\ex \textbf{Το πλάγιο.} Ένα σημείο $\mathrm{Α}$ απέχει $4$ cm από μια ευθεία $(\varepsilon)$. Ένα σημείο $\mathrm{Β}$ της $(\varepsilon)$ μπορεί να απέχει από το $\mathrm{Α}$ $3$ cm; $4$ cm; $6$ cm;

\ex \textbf{Το ράφι.} Ένα ράφι είναι παράλληλο προς το πάτωμα και απέχει $1{,}2$ μ. Πόσο απέχει κάθε σημείο του από το πάτωμα;

\ex \textbf{Το χωράφι.} Ένα ορθογώνιο χωράφι έχει διαστάσεις $40$ και $25$ μ. Πόση είναι η απόσταση ανάμεσα στις μεγάλες απέναντι πλευρές του; Και ανάμεσα στις μικρές;

\zoni{Το σκέφτομαι}{ο συντομότερος δρόμος}

\ex Εξήγησε γιατί το κάθετο τμήμα είναι το συντομότερο (σκέψου ένα τεντωμένο νήμα και μια ράβδο που στρέφεται).

\ex Σε ποια περίπτωση η απόσταση δύο παραλλήλων είναι $0$;

\ex Βρες τα σημεία που απέχουν $3$ cm από την $(\varepsilon)$ και $5$ cm από μια παράλληλη προς αυτήν σε απόσταση $8$ cm.

\ex Δύο παράλληλες ευθείες απέχουν $6$ cm. Πόσες ευθείες παράλληλες προς αυτές απέχουν $3$ cm και από τις δύο;

\zoni{Για το σπίτι}{μέτρησε κάθετα}

\ex Μέτρησε την απόσταση ανάμεσα σε δύο παράλληλες ακμές του τραπεζιού σου σε δύο διαφορετικά σημεία.

\ex Βρες τη μικρότερη απόσταση ενός σημείου του δωματίου από έναν τοίχο και εξήγησε τι μέτρησες.
